\documentclass[10pt,a4paper]{amsart}
\usepackage[margin=19mm]{geometry}
\usepackage{amsmath,amssymb,mathtools}
\usepackage[scr=rsfs,cal=euler]{mathalfa}
\usepackage{xfrac}
\usepackage[unicode,hidelinks,bookmarks=false]{hyperref}
\newcommand{\ie}{i.e.\ }
\newcommand{\cf}{cf.\ }
\newcommand{\resp}{respectively\ }
\newcommand{\Subsection}{\subsection}
\providecommand{\nobreakdash}{\nobreak\hbox{-}}
\setlength{\emergencystretch}{2em}
\multlinegap0pt
\usepackage{bm}
\usepackage{bbm}
\usepackage{dsfont}
\usepackage{xspace}
\usepackage{cancel}
\usepackage{mathtools}
\usepackage{amsthm}
\makeatletter
\@ifundefined{prop}{\newtheorem{prop}{Proposition}}{}
\makeatother
\makeatletter
\expandafter\ifx\csname red\endcsname\relax
\newenvironment{red}{\relax\color{red}}{\hspace*{.5ex}\relax}
\fi
\makeatother
\title[From Group Operations to Geometric Structures: Amalgamations]{From Group Operations to Geometric Structures: Amalgamations, HNN-Extensions, and Twisting in Coset Geometries}
\author{Claudio Alexandre Piedade, Philippe Tranchida}
\date{Source: arXiv:2505.24662v1 (CC BY 4.0)}
\begin{document}
\maketitle

\noindent\emph{Source and licence.} From Group Operations to Geometric Structures: Amalgamations, HNN-Extensions, and Twisting in Coset Geometries. Version arXiv:2505.24662v1 (CC BY 4.0), distributed under the Creative Commons Attribution 4.0 International licence, as recorded in the arXiv licence metadata for this exact version and shown on the abstract page https://arxiv.org/abs/2505.24662v1. Full rights notice: licenses/arxiv-2505.24662-CC-BY-4.0.txt.\par\smallskip

\noindent\textit{Editorial scope.} Counted unit: the Prop whose source label is \texttt{prop:semidir\_FT} in the article (source0.tex lines 1162--1164, source label \texttt{prop:semidir\_FT}), delivered together with the article's own complete original proof of it (source0.tex lines 1165--1201, one \texttt{proof} environment of 37 lines). No mathematics is written, repaired or re-derived: statement and proof are the article's own text line for line. Required context retained uncounted: none. Every definition, display and label of those lines is retained unchanged. Omitted from this package and not redistributed: 1--1161, 1202--2006. Nothing inside a retained excerpt is omitted or altered: every retained body is delivered exactly as the article's lines are written, so no omission receipt of any kind accompanies this package. Citations: the retained excerpts carry no citation command at all, so no bibliography entry of the article is transcribed and no reference is redistributed. Reference adaptation: none was needed and none was made; every reference inside the delivered unit resolves inside it, so no unresolved reference or citation is produced. Printed slips kept literal: no printed word, symbol, hypothesis or number of the counted statement or of its proof was altered. Layout and class: the article's own class is \texttt{amsart} with packages including inputenc, amsmath, amssymb, graphicx, setspace, verbatim, url, color, hyperref, multicol, enumitem, float, soul, tikz; the wrapper is the lane's pinned \texttt{amsart} editorial wrapper at 10pt on A4, which restates the theorem-like environments the retained text needs and adds the article's own macro definitions that the retained text needs (none), with extra wrapper packages (bm, bbm, dsfont, xspace, cancel, mathtools). The delivered numbering is the unit's own. Assets: no retained span contains a figure environment, a graphics-inclusion command, a PDF-page inclusion, a table or a TikZ picture; no image file is copied into this package, the package declares no asset dependency and no image file is redistributed with it. Licence: version arXiv:2505.24662v1 of this article is distributed under the Creative Commons Attribution 4.0 International licence, as recorded in the arXiv licence metadata for this exact version and shown on the abstract page https://arxiv.org/abs/2505.24662v1. A full rights notice is delivered at licenses/arxiv-2505.24662-CC-BY-4.0.txt. Characters: every retained excerpt is pure ASCII, so the delivered engine, XeLaTeX with its default Latin Modern text font, reports no missing character.
\begin{prop}\label{prop:semidir_FT}
If $A$ is flag-transitive on $\alpha = (A,(A_i)_{i \in I_\alpha})$ and $B$ is flag-transitive on $\beta = (B,(B_i)_{i \in I_\beta})$, then $G = A \rtimes_\varphi B$ is flag-transitive on $\alpha \rtimes_\varphi \beta$.
\end{prop}

\begin{proof}
    We want to show that $G_JG_i = \cap_{j\in J} G_j G_i$ for any $J \subseteq I$ and $i\in I\setminus J$. Let $J = J_\alpha \sqcup J_\beta$, with $J_\beta \subseteq I_\beta$ and $J_\alpha \subseteq K$.
    
    First suppose that $i \in I_\beta$. Then, the equation we need to verify becomes

    $$ (A_{J_\alpha} \rtimes B_{J_\beta}) (A \rtimes B_i) = (\cap_{j \in I_\beta \cap J} (A \rtimes B_j) (A \rtimes B_i)) \cap ( \cap_{k \in K \cap J} (A_j \rtimes B) (A \rtimes B_i)) $$

    We will once again work with the underlying sets of every subgroups appearing in the equation. In doing so, the equation splits into the following two equations.
    
    \begin{equation}\label{eq:SDFT1}
          A_{J_\alpha} A^{B_{J_\beta}} =  \cap_{k \in K \cap J} A_j A^{B_i}
    \end{equation}
    where $A^{B_{J_\beta}}$ denotes the set of all elements of the form $a^{b}$ with $a \in A$ and $b \in B_{J_\beta}$,
    and 
  \begin{equation}\label{eq:SDFT2}
      B_{J_\beta}B_i = \cap_{j\in I_\beta \cap J} B_j B_i.
  \end{equation}
    Equation \ref{eq:SDFT2} holds directly from the fact that $B$ acts flag-transitively on $\beta$. For Equation \ref{eq:SDFT1}, note that $A^{B_i} = A^{B_{J_\beta}} = A$, so the equation holds trivially.

    Suppose now that $i \in K$. The equation then becomes

        $$ (A_{J_\alpha} \rtimes B_{J_\beta}) (A_i \rtimes B) = (\cap_{j \in I_\beta \cap J} (A \rtimes B_j) (A_i \rtimes B) )\cap (\cap_{k \in K \cap J} (A_k \rtimes B) (A_i \rtimes B))$$


which then splits into

\begin{equation}\label{eq:SDFT3}
          A_{J_\alpha} A_i^{B_{J_\beta}} =  \cap_{k \in K \cap J} A_k A_i^{B}
    \end{equation}
    and 
  \begin{equation}\label{eq:SDFT4}
      B_{J_\beta}B = \cap_{j\in I_\beta \cap J} B_j B
  \end{equation}

  This time, Equation \ref{eq:SDFT4} holds trivially. For Equation \ref{eq:SDFT3}, note that since $i\in K$, $i$ is a orbit of types closed under the action $\varphi$. Hence, $A_i^{B_{J_\beta}} = A_i =A_i^B$, so that the equation holds by flag-transitivity of $A$ of $\alpha$.

\end{proof}

\end{document}
